% !PRINTED 485
\weberpage{485}

eine ganze rationale Function von $t$, deren Co\"efficienten rationale
Functionen der $\alpha, \alpha_1 \ldots \alpha_{n-1}$ sind, die ungeändert bleiben,
wenn irgend eine Permutation der Gruppe $P$ vorgenommen wird,
weil durch diese Permutationen die $\omega$ entweder ungeändert
bleiben, oder in derselben Weise wie die $\Theta$ mit einander permutirt werden. Nach \S.\,149 b) sind diese Co\"efficienten also in $\Omega$
enthalten. Setzen wir dann in (5) für das unbestimmte $t$ den
Werth $\Theta$ ein, so folgt, da $\varphi'(\Theta)$ von Null verschieden ist,
\weq{6}{\omega = \frac{\Phi(\Theta)}{\varphi'(\Theta)}.}

Wenden wir dies an auf die Co\"efficienten des Productes
\weq{7}{(u - \alpha)(u - \alpha_1) \ldots (u - \alpha_{r-1}) = \varphi(u, \Theta),}
wo $u$ eine Variable bedeutet, so ergiebt sich, dass diese Function
$r^{\text{ten}}$ Grades, deren Wurzeln die $\alpha, \alpha_1 \ldots \alpha_{r-1}$ sind, rational
durch $\Theta$ ausgedrückt werden kann.

Es ist also die Grösse $\alpha$, die ursprünglich Wurzel einer
Gleichung $n^{\text{ten}}$ Grades in $\Omega$ ist, zugleich Wurzel einer Gleichung
$r^{\text{ten}}$ Grades, deren Co\"efficienten rational von der Wurzel einer
Gleichung $s^{\text{ten}}$ Grades abhängen. Eine Gleichung $f(x) = 0$,
deren Wurzel $\alpha$ diese Eigenschaft hat, nennt man eine \gesperrt{imprimitive Gleichung}.

Wir können das Bewiesene auch so ausdrücken:

\gesperrt{2. Der imprimitive Körper $\Omega(\alpha)$ vom $n^{\text{ten}}$ Grade geht
durch Adjunction des Körpers $s^{\text{ten}}$ Grades $\Omega' = \Omega(\Theta)$
zu $\Omega$ in einen Körper $\Omega'(\alpha)$ vom $r^{\text{ten}}$ Grade über.}

Wir wollen noch untersuchen, ob die Imprimitivität der
Gruppe ein ausreichendes Kennzeichen für die Imprimitivität des
Körpers ist, ob man also in einem Körper $\Omega(\alpha)$ mit imprimitiver
Gruppe immer imprimitive Elemente finden kann.

Sei also jetzt $f(x) = 0$ eine irreducible Gleichung $n^{\text{ten}}$
Grades mit imprimitiver Gruppe, und seien (1) die Wurzeln
dieser Gleichung, die so in die Reihen (2) zerlegt sind, dass die
Elemente dieser Reihen durch die Permutationen der Gruppe
nicht von einander getrennt werden.

Wir wählen irgend eine symmetrische Function $\psi(x, x_1 \ldots x_{r-1})$
so, dass die Werthe
%CONT>>
